\documentclass[12pt,oneside]{report}

% ---------------------------------------------------------------------------
% Governed Autonomy: A Manufactured Phase Change
% A doctoral thesis manuscript grounded in the AutoFDE Lab artifact chain.
%
% Compile: latexmk -pdf phase-change-thesis.tex
%
% Every quantitative claim in this manuscript is anchored to a measured
% artifact; Appendix A is the verification map. Negative results are
% first-class citizens (Chapter 6): the claims that did not survive
% adversarial audit are cited by the commit that deleted them.
% ---------------------------------------------------------------------------

\usepackage[T1]{fontenc}
\usepackage[utf8]{inputenc}
\usepackage{lmodern}
\usepackage[margin=1.1in]{geometry}
\usepackage{microtype}
\usepackage{amsmath}
\usepackage{amssymb}
\usepackage{amsthm}
\usepackage{booktabs}
\usepackage{listings}
\usepackage[hidelinks]{hyperref}

\lstset{basicstyle=\ttfamily\small, breaklines=true, frame=single,
        columns=fullflexible, keepspaces=true}

\newtheorem{invariant}{Invariant}
\newtheorem{principle}{Principle}
\newtheorem{definitionx}{Definition}

\newcommand{\code}[1]{\texttt{#1}}
\newcommand{\commit}[1]{\texttt{#1}}
\newcommand{\daggerfootnote}{$\dagger$\ }

\title{\textbf{Governed Autonomy as a Manufactured Artifact}\\[6pt]
\large The Deletion of Implementation as a Phase Transition\\
\large in the Engineering of Autonomous Software Agents}
\author{Sean Chatman\\
\small AutoFDE Laboratory\\
\small \texttt{sean@chatmangpt.com}}
\date{September 2026}

\begin{document}

\maketitle

\begin{abstract}
\noindent
Autonomous software agents must allocate finite exploration resources under
uncertainty, and every existing family of allocation policies---greedy
selection, $\varepsilon$-greedy exploration, confidence bounds (UCB1), and
predictor--confidence upper-confidence trees (PUCT)---optimizes expected
performance while remaining structurally incapable of governance: their
outputs are non-deterministic across platforms, non-conserving under
discretization, silently degrading under failure, and invisible to
after-the-fact audit. This thesis argues and demonstrates that the binding
constraint on trustworthy agent autonomy is not exploration performance but
the \emph{auditability of consequential decisions}, and that removing this
constraint \emph{structurally}---rather than mitigating it---constitutes a
technological phase change in the sense established by the steam engine and
the transistor: each removed an entire failure class of its predecessor,
and each thereby unlocked composition (the factory system; integrated
circuits) that the old regime could not reach at any price.

The demonstrated artifact is a single certified allocation kernel
(\code{bcinr-cmca}, branchless, cyclomatic complexity $CC=1$, Q16.16
fixed-point) that serves every host in a polyglot ecosystem bit-identically:
in-process through WebAssembly on both CPython (via \code{wasmtime}) and the
BEAM virtual machine (via \code{wasmex}), with agreement pinned by golden
fixtures whose values are exact binary fractions. The work's methodological
core is a \emph{falsification ledger}: the thesis deliberately reports and
cites the deletion of its own weaker claims---a circular comparative
benchmark, a set of ``theorems'' that held for any softmax, a notational
equivalence proof, and ultimately the entire local reference implementation
of the allocation measure (removed in commit \commit{c1d04ffc}). What
survives falsification is the governance envelope: deterministic replay,
conservation by construction, typed refusal, and content-bound receipts
streamed into Object-Centric Event Log (OCEL~2.0) models where they are
conformance-checked and counterfactually mutated.

The empirical record includes an honest negative finding: the certified
kernel's compiled lens policy is substantially \emph{flatter} than the
deleted tunable softmax (a dead-end candidate receives a measured $16.9\%$
share where the softmax at high inverse temperature gave it none), a
trade-off we document rather than tune away. The phase-change claim is
argued structurally---failure-class removal, composition unlock, and a
compounding standard interface---and scoped honestly: this is a
laboratory demonstration of structure, not a demonstrated historical
shift. Falsifiers for the analogy are stated explicitly.
\end{abstract}

\tableofcontents

\chapter{Introduction}
\label{ch:intro}

\section{The problem that would not go away}

Software agents that act on the world must decide, repeatedly and under
uncertainty, how to divide finite resources---compute ticks, memory,
verification depth, consequence budget---across candidate courses of
action. The classical answers are well understood and theoretically
characterized: greedy exploitation, $\varepsilon$-greedy exploration,
UCB1's optimism in the face of uncertainty~\cite{auer2002}, and PUCT-style
search~\cite{rosin2011}. Each is a performance instrument. None is a
governance instrument. Their outputs cannot in general be replayed
bit-identically across platforms; their budgets leak under discretization;
their failures degrade silently through fallbacks; and their decisions
leave no durable, checkable trace.

For toy systems this does not matter. For \emph{consequential} autonomy---
agents whose actions carry authority, cost, or risk---it is the binding
constraint. An unauditable allocation is an unaccountable one, and an
unaccountable allocation cannot be delegated, chained, or composed across
trust boundaries. The hypothesis of this thesis is that this constraint,
not exploration suboptimality, is what currently prevents the composition
of autonomous agents into large, trustworthy systems.

\section{The thesis statement}

\begin{principle}[Thesis statement]
\label{pr:thesis}
When the allocator of an autonomous agent's consequential resources is
manufactured as a single certified artifact---deterministic across hosts,
conserving by construction, refusing by type rather than degrading
silently, and emitting content-bound receipts into a process-mining event
log---the failure class of \emph{divergent, silent, unverifiable
implementation} is removed structurally. This removal unlocks composition
(polynomial numbers of hosts, long chains of receipted delegation,
conformance-checked lifecycles) that per-implementation governance cannot
reach at any tuning budget, in the same structural sense in which the
steam engine removed the biological-power failure class and the transistor
removed the filament failure class.
\end{principle}

Three clarifications keep this claim honest. First, the analogy is
\emph{structural}, not historical: we demonstrate a mechanism, not a
societal shift. Second, at the moment of transition the new regime may be
\emph{worse} on legacy metrics---Newcomen's engine was thermally
inefficient, the first transistors were not faster than tubes, and our
certified kernel is flatter than a tuned softmax; Chapter~\ref{ch:phase}
argues that this is characteristic, not incidental. Third, the claim is
falsifiable, and Chapter~\ref{ch:phase} states what would falsify it.

\section{Contributions}

\begin{enumerate}
\item \textbf{The governance envelope, formalized}
(Chapter~\ref{ch:envelope}): a precise statement of the four laws a
consequential allocator must satisfy---deterministic replay, conservation,
typed refusal, receipt---with an explicit demotion of the allocation
\emph{mathematics} to standard constructions (escort measures over
multifractal cascades) and an account of why the envelope, not the
measure, carries the contribution.
\item \textbf{The single-kernel law} (Chapter~\ref{ch:kernel}): an
architecture and a verified deployment in which one certified Rust kernel
is the only implementation of the allocation measure in a polyglot
ecosystem, served in-process via WebAssembly to CPython and the BEAM,
with bit-identical outputs pinned by cross-host golden fixtures. The
local reference implementation was deleted outright, and its deletion is
pinned by a test that makes reintroduction a structural type error.
\item \textbf{Process-law conformance for allocations}
(Chapter~\ref{ch:process}): allocation plans as first-class OCEL~2.0
events, checked by alignments against discovered process models and
adversarially mutated by a counterfactual battery.
\item \textbf{The falsification ledger as method}
(Chapter~\ref{ch:ledger}): a documented, commit-cited record of claims
made, claims falsified, and claims survived---including the autopsy of
our own circular benchmark. We propose deletion-with-receipt as a
first-class scholarly act.
\item \textbf{A phase-change analysis} (Chapter~\ref{ch:phase}): the
constraint-removal model of technological phase transitions, applied with
explicit disanalogies and falsifiers.
\end{enumerate}

\section{Why this needed a laboratory, not just a paper}

Every claim above lives or dies on artifacts: a kernel that exists, hosts
that agree bit-for-bit, tests that fail when the law is broken, and a
commit history in which the negative results are visible. The thesis is
therefore written \emph{from} the artifact chain summarized in
Appendix~\ref{app:artifacts}; where a number appears in the text, it was
measured, and the measurement is reproducible from the cited commit.

\chapter{Background and Related Work}
\label{ch:background}

\section{Allocation under uncertainty}

The multi-armed bandit literature establishes the performance frontier for
sequential allocation: UCB1's logarithmic regret~\cite{auer2002}, PUCT's
confidence-driven search~\cite{rosin2011}, and the general treatment in
Sutton and Barto~\cite{sutton2018}. This literature optimizes reward under
a trust model in which the \emph{implementor} of the policy is trusted
absolutely. The present work assumes the opposite trust model---the
policy's decisions must be verifiable by parties who did not implement
it---and therefore asks a question the bandit literature does not pose:
what is the \emph{decision artifact} that a third party can replay,
conserve, and audit?

\section{Deterministic arithmetic and certified kernels}

Floating-point divergence across platforms is classical,
well-documented~\cite{goldberg1991}, and fatal to bit-identical replay.
The standard remedy is fixed-point arithmetic; the kernel at the center of
this thesis uses Q16.16 (sixteen integer bits, sixteen fractional bits,
resolution $2^{-16}$) with saturating operations whose truncation is
one-sided. The deeper tradition is certified code: kernels constrained to
cyclomatic complexity $CC=1$~--- straight-line, branchless,
allocation-free---inherit constant-time execution and side-channel
resilience as theorems of their shape rather than properties of their
testing. MacKenzie's study of mechanized proof~\cite{mackenzie2001}
frames the sociological point this thesis operationalizes: trust in
computation migrates from persons to artifacts only when the artifacts
carry their own evidence.

\section{Object-centric process mining}

The Object-Centric Event Log standard (OCEL~2.0)~\cite{ocel2,pm4py}
records events tied to multiple objects with qualifiers, enabling process
conformance over entities rather than flattened traces. Object-centric
process querying extends this to declarative structural laws over event
logs~\cite{ocpq}. The present work uses OCEL~2.0 as the \emph{durably
checkable substrate} for governance: an allocation that becomes an event
with a content-bound hash can be conformance-checked, counterfactually
mutated, and replayed long after the process that produced it is gone.

\section{Hierarchical planning and agent execution}

FOND planning and HDDL hierarchical task networks provide the exploration
structures over which allocation happens; GymAct-style governed
simulation provides execution. These are the consumers of the allocator
and are treated as given, cited in Appendix~\ref{app:artifacts} as
artifact rather than surveyed here as literature.

\section{The industrial analogues}

Landes' account of the industrial revolution~\cite{landes1969},
Rosenberg's studies of technological change~\cite{rosenberg1976}, and
Mokyr's economic history of invention~\cite{mokyr1990} establish the
pattern this thesis borrows: general-purpose technologies transform when a
\emph{constraint class} (fuel availability at the minehead, component
reliability) is removed and improvement effort compounds onto the new
regime. The transistor's own record---Bardeen and
Brattain's 1948 demonstration~\cite{bardeen1948} and Moore's 1965
composition argument~\cite{moore1965}---supplies the cleanest instance:
the initial win was reliability (no filament), and the compounding came
from composition (integration). Chapter~\ref{ch:phase} develops this.

\chapter{The Governance Envelope}
\label{ch:envelope}

\section{What CMCA is, honestly}

Chatman Multifractal Cascade Allocation (CMCA) names the following
construction. Given candidate branches $c_1,\dots,c_N$ with option entropy
$\Omega(c_i)$, historical yield $Y(c_i)$, and cost $C(c_i)$, define the
salience density
\begin{equation}
S(c_i) \;=\; \frac{\Omega(c_i)\, Y(c_i)}{\max(C(c_i), \varepsilon)},
\label{eq:salience}
\end{equation}
and allocate mass by a multifractal escort cascade: at each parent with
children, the local escort distribution at lens $q$ is
\begin{equation}
L_q(c_i \mid p) \;=\; \frac{m(c_i)^q}{\sum_j m(c_j)^q},
\qquad m(c) = e^{S(c)},
\label{eq:escort}
\end{equation}
and a node's consequence mass is the product of local escorts down its
root path. Setting $m = e^{S}$ makes~\eqref{eq:escort} a
Gibbs--Boltzmann softmax at inverse temperature $\tau = q$; the
``multifractal'' content is the per-level lens vector $(q_0, q_1, \dots)$,
which the certified kernel compiles from an ontology-declared policy.

We state plainly, as a matter of scholarly hygiene and as the first entry
of the falsification ledger (Chapter~\ref{ch:ledger}): \emph{the
mathematics is standard}. Equation~\eqref{eq:salience} is a heuristic
payoff density; equation~\eqref{eq:escort} is the textbook escort
distribution of multifractal analysis; the one-line identity
$dH/d\tau = -\tau\,\mathrm{Var}(S) \le 0$ for the softmax's entropy is an
exponential-family fact. An earlier iteration of this work presented
these as ``theorems'' of CMCA; audit showed they hold for \emph{any}
softmax allocator, and the corresponding test file was deleted
(commit-recorded in Appendix~\ref{app:artifacts}). What survives audit,
and what this chapter exists to formalize, is the \emph{envelope}---the
four laws the artifact satisfies that no stock exploration policy does.

\section{The four laws}

\begin{invariant}[Deterministic replay]
For identical inputs, the allocator produces bit-identical outputs on
every host, and the output carries a canonical hash: candidates are
ordered by a canonical \code{candidate\_hash} before any arithmetic, and
all arithmetic is integer fixed-point. Inspection is not enough; the
cross-host golden fixture of Chapter~\ref{ch:kernel} is the executed
proof.
\end{invariant}

\begin{invariant}[Conservation by construction]
Discrete budgets are assigned by floor projection:
$\sum_i \lfloor p_i B \rfloor \le B \sum_i p_i = B$, with the residual
one-sided and bounded by the number of active branches. The certified
kernel's own documentation goes further than we initially did: it
\emph{measured} its residuals across adversarial fixtures before claiming
the bound, and explicitly narrowed an earlier ``leaves sum to one''
claim after finding a legitimate exception. Conservation is a property of
the projection, verified on prime budgets (e.g., $104{,}726 \le 104{,}729$
ticks over an 8-branch frontier, residual within the bound).
\end{invariant}

\begin{invariant}[Typed refusal]
Every failure mode refuses with a named type rather than degrading:
cardinality beyond the compiled $N{=}8$ shape, non-finite features, a
forced-but-absent transport, a mis-pointed artifact path, numeric
overflow, underflow of an escort weight under a negative lens. The
canonical failure this law exists to prevent is the \emph{silent
fallback}: during this work, a test caught our own bridge silently
falling through a mis-set environment variable to a different artifact;
the law was tightened and the test pinned (the refusal now names the
offending path).
\end{invariant}

\begin{invariant}[Receipt]
Every allocation plan is content-bound: a SHA-256 digest over a canonical
projection of the plan. Receipts stream into OCEL~2.0 event logs, where
they become ordinary process data---queryable, conformance-checkable,
and counterfactually mutable (Chapter~\ref{ch:process}).
\end{invariant}

\section{What the envelope buys, and what it costs}

The envelope buys composition: any host that can execute WebAssembly can
serve the identical law, and any auditor with the receipt can replay the
decision. It costs performance flexibility: there is no temperature knob.
The compiled lens policy is what it is; changing exploration behavior
means regenerating policy upstream (a governance act with provenance),
not tuning a parameter at the call site. The empirical consequence is
measured in Chapter~\ref{ch:kernel}: the compiled policy is
substantially flatter than a tuned softmax. We regard this as the
transistor-versus-filament trade in miniature---a capability exchanged
for a guarantee---and we report it rather than tune it away.

\chapter{The Single-Kernel Law}
\label{ch:kernel}

\section{From one engine per host to one kernel for all hosts}

Before this work, the ecosystem contained a certified Rust kernel
(\code{bcinr-cmca}), a subprocess CLI wrapping it, and---as a local
convenience---a float-softmax reference implementation inside the Python
package. The reference implementation was the standard failure waiting to
happen: a second implementation of the measure that would drift, diverge
in its floats, and eventually be trusted by someone who did not know it
was a stand-in. The single-kernel law is the deletion:

\begin{quote}
\emph{There is exactly one implementation of the allocation measure in
the ecosystem. It is the vendored, certified kernel. Every host obtains
its measure from that kernel. There is no fallback implementation,
reference or otherwise.}
\end{quote}

Enforcement is structural, not conventional: the Python allocator class
no longer accepts an \code{engine} or \code{tau} argument at all (a
Chicago test pins that reintroducing them raises \code{TypeError}), and
the surviving transports---in-process WebAssembly first, CLI subprocess
as availability fallback---are built from the same vendored commit
(\commit{2ae60cf0}), so transport choice cannot change an answer.

\section{The WebAssembly deployment}

The kernel is compiled to \code{wasm32-wasip1} as a cdylib exporting a
three-function linear-memory ABI:

\begin{lstlisting}
cmca_alloc(len) -> ptr
cmca_free(ptr, len)
cmca_rank_json(ptr, len) -> ptr    ; [u32 LE out_len][JSON bytes]
\end{lstlisting}

\noindent The request/response JSON contract is the CLI's, verbatim:
$\{$``candidates'': $[\{$``name'', ``measures''$[m_0..m_3]\}]$\}$ in;
$\{$``ranking'': $[\{$``name'', ``share''$\}]$\}$ out; refusals as
$\{$``error''$\}$ envelopes with the engine's own message. The four
caller-defined measure axes are the branch's option entropy, inverse
cost, historical yield, and salience~\eqref{eq:salience}---the same
feature vector for every host, encoded once, in the bridge.

Two hosts were verified. CPython loads the artifact through
\code{wasmtime}; the BEAM loads the \emph{same artifact file} through
\code{wasmex}, with trap-restart semantics (the guest builds with
\code{panic=abort}, so a trapped instance is killed and rebuilt, its
state never reused). The artifact is committed to the repository
(421{,}988 bytes), so consumers need neither a Rust toolchain nor a
subprocess binary---only a WASM runtime.

\section{Bit-identity across hosts}

The load-bearing empirical claim of this thesis is this table. One
request---three candidates with fixed measure axes---was evaluated by the
kernel through four paths: natively (CLI subprocess), in-process on
CPython (wasmtime), in-process on the BEAM (wasmex), and the raw JSON
compared across transports:

\begin{table}[h]
\centering
\begin{tabular}{lll}
\toprule
Candidate & Share (exact) & As binary fraction \\
\midrule
beta  & $0.2135009765625$ & $13{,}983/65{,}536$ \\
alpha & $0.153564453125$  & $10{,}062/65{,}536$ \\
gamma & $0.1290130615234375$ & $8{,}452/65{,}536$ \\
\bottomrule
\end{tabular}
\caption{Cross-host golden fixture. The shares are exact Q16.16 binary
fractions; the BEAM test pins these values, captured through the
CPython/wasmtime path over the same artifact. Any drift in the artifact,
the transport, or the host binding fails the fixture by name.}
\label{tab:golden}
\end{table}

The payloads are byte-identical JSON across all four paths. This is not
an approximation claim to be argued; it is an equality claim to be
tested, and the test is in the repository. To our knowledge, the
deployment of one certified fixed-point policy kernel serving Python and
Erlang hosts bit-identically through WebAssembly, with the agreement
pinned by golden fixtures, is a construction not previously documented in
the agent-governance literature.

\section{The honest finding: flatness}

When the tunable softmax was deleted, the benchmarks that swept its
inverse temperature were refactored to drive the shipped kernel. The
measured behavior of the compiled lens policy is markedly flatter than a
high-$\tau$ softmax. In a scenario constructed to contain one viable
pair and one dead-end candidate (entropy $0.01$, yield $0.01$, cost
$100$), the compiled policy allocates:

\begin{table}[h]
\centering
\begin{tabular}{lc}
\toprule
Candidate & Measured share \\
\midrule
b\_promising     & $0.490498$ \\
b\_promising\_alt & $0.340867$ \\
b\_dead\_end     & $0.168635$ \\
\bottomrule
\end{tabular}
\caption{The compiled lens policy funds every candidate with real mass;
the deleted softmax at $\tau{=}5$ starved the dead end entirely. The
pruning cutoff is therefore a projection-level parameter calibrated to
measured shares, documented in the test that pins it.}
\label{tab:flat}
\end{table}

The interpretation offered to the reader is deliberately double-edged.
As a performance instrument, the compiled policy leaves exploration
budget on the table by the standards of a tuned softmax. As a governance
instrument, its flatness is honest: the policy refuses to silently
murder a candidate someone lawfully proposed, and forces the kill
decision into the projection layer where it is receipted. Which of these
framings matters is precisely the question the thesis exists to re-pose.

\chapter{Process Law and Counterfactual Conformance}
\label{ch:process}

\section{Allocations as events}

An allocation plan that ends its life in application memory is
ungoverned. The receipted plan is projected into OCEL~2.0 logs as events
bound to the objects they concern (plans, branches, worlds), alongside
the actuation receipts of the actions the allocation authorized. From
that moment the allocation is process data: queryable with
object-centric query languages, joinable across the lifecycles of
interacting agents, and discoverable as process models.

\section{Conformance and counterfactual mutation}

Conformance checking against discovered process models
(alignments~\cite{pm4py}) validates that executed lifecycles obey the
process law. The counterfactual battery goes further: programmatic
mutants---premature actuation, unadmitted action injection, skipped
teardown, subtask inversion, crossed object links, resource-allocation
inversion---must each be \emph{detected} by the governed pipeline, and
the detection is tested. The resource-allocation mutant is the one this
thesis adds to the literature's usual set: it inverts the allocator's own
output (assigning the strongest branch the weakest budget) and requires
that the inversion fail conformance downstream. A governor whose
decisions cannot be counterfactually falsified is decoration; the battery
is the difference.

\section{The limitation we declare}

The mutators are generated by the same codebase that generates the
defenses. Known-mutation testing is strong regression discipline but
weak adversarial evidence; it certifies the fence against the adversaries
the fence-builder imagined. Chapter~\ref{ch:ledger} records this as an
open standing item rather than a resolved one.

\chapter{The Falsification Ledger}
\label{ch:ledger}

\section{Deletion as a scholarly act}

A thesis that only accumulates claims trains its author to defend the
indefensible. This manuscript's methodological core is the opposite
discipline: every claim is adversarially audited, and claims that fail
are \emph{deleted with receipt}---removed from the codebase in named
commits, recorded here, and cited as evidence that the audit was real.
The ledger also enforces a monotone law borrowed from the laboratory's
operating contract: the hand-written surface must shrink; growth of the
hand-written surface requires, in the same change, a paydown plan.

\section{The autopsies}

\subsection{The circular benchmark}

An earlier ``empirical value analysis'' compared CMCA against a greedy
baseline defined as argmax on immediate yield---an allocator blind, by
construction, to the option-entropy term in the evaluation metric, which
was in turn the allocator's own objective. The scenario was hand-tuned so
the entropy-aware algorithm had to win; the comparator's arm existed only
as comments. Audit verdict: a characterization test masquerading as a
comparative experiment. Deleted as evidence; retained (rewritten) as a
characterization of the shipped engine, with its thresholds calibrated to
measured shares and its numbers documented in-test
(Table~\ref{tab:flat}).

\subsection{The theorems that held for every softmax}

Four ``theorems'' pinned properties of the softmax measure: floor
conservation ($\sum\lfloor p_iB\rfloor \le B$, one line from the
definition of floor), entropy monotonicity ($dH/d\tau = -\tau
\mathrm{Var}(S)$, an exponential-family identity), order preservation
(softmax monotonicity), and permutation invariance (sorting commutes).
Audit verdict: true, elementary, and non-discriminating---they
distinguished CMCA from no softmax allocator whatsoever. The test file
was deleted with the reference engine (commit \commit{c1d04ffc}).

\subsection{The notational equivalence}

A formal ``proof of equivalence'' between the Python softmax and the
kernel's escort distribution reduced, on inspection, to the observation
that $m = e^S$ turns $m^q/\sum m^q$ into $e^{qS}/\sum e^{qS}$: a change
of notation, verified numerically to $10^{-6}$, which confirmed only that
floating-point arithmetic worked. Deleted. The equivalence that matters
was rebuilt as the bit-identity fixture of Table~\ref{tab:golden}, which
compares real builds of real transports rather than formulas with
themselves.

\subsection{The reference implementation itself}

The final and largest entry: the local float-softmax engine, its
temperature parameter, its salience method, and its end-to-end plumbing
(plan fields, port contracts, cross-language resource attributes) were
deleted outright once the single-kernel law held. The deletion removed
approximately a thousand lines of hand-written measure code and every
test that existed only to pin its behavior. What remains in the Python
package is a projection: ordering, pruning, renormalization,
anti-starvation, discrete assignment---all engine-agnostic, all tested
against the kernel.

\section{What the ledger buys}

The ledger converts the thesis's central risk---that a governance
claim is itself an unaudited assertion---into structure. The reader
does not have to believe the artifact is honest; the reader can observe
that dishonesty would have required \emph{surviving} an audit process
that demonstrably deleted four of its own claims, including its own
reference implementation. The ledger is, in the strict sense, the
thesis's own governance envelope applied to itself.

\chapter{The Phase-Change Analysis}
\label{ch:phase}

\section{A model of phase change}

We adopt the following structural model, synthesized from the economic
history of general-purpose technologies~\cite{landes1969,rosenberg1976,
mokyr1990}. A technology undergoes a phase change when three conditions
hold jointly:

\begin{enumerate}
\item \textbf{Constraint-class removal.} An entire \emph{class} of
failures or limitations of the old regime is removed structurally---not
reduced, not mitigated, but made impossible by construction.
\item \textbf{Composition unlock.} The removal enables assemblies of a
kind and scale that the old regime could not reach at any price,
because the old failure rate grew with the size of the assembly.
\item \textbf{A compounding interface.} A standardized interface causes
improvement effort to accumulate onto the new regime rather than
forking across variants.
\end{enumerate}

\subsection{Steam}

The steam engine did not outrun good water mills on efficiency at
introduction; early engines were thermally poor and economical only
where fuel was free (minehead pumping). The phase change was the removal
of the \emph{biological and geographical} constraint class: power that
did not tire, did not freeze, and did not depend on river siting. The
composition unlock followed---factories could locate anywhere, and the
factory system emerged as an assembly of powered stations. The
compounding interface was standardization (interchangeable parts,
standard threads), which turned power into a commodity input that
improved system-wide when any engine improved.

\subsection{The transistor}

At demonstration in 1947--48~\cite{bardeen1948}, the transistor was not
faster than the vacuum tube. Its initial superiority was the removal of
the \emph{filament failure class}: no thermal wear-out, no warm-up,
lower power. The composition unlock is the heart of the story: a machine
with $n$ components of per-component reliability $r$ has availability
$r^n$, and ENIAC-scale tube assemblies were already failing continuously;
million-component machines were arithmetically impossible in tubes. The
removed failure class is what made large assemblies thinkable, and
integration~\cite{moore1965} is the compounding interface---the planar
process---onto which fifty years of improvement effort accumulated.

\section{The mapping}

\begin{table}[h]
\centering
\small
\begin{tabular}{p{0.28\textwidth}p{0.32\textwidth}p{0.32\textwidth}}
\toprule
 & \textbf{Old regime: per-implementation policies} & \textbf{This work: single-kernel governance} \\
\midrule
Constraint class removed &
Divergent implementations; float non-portability; silent fallbacks; invisible decisions &
One certified kernel; integer fixed-point replay; typed refusal; receipted events \\
Composition unlock &
Audit surface grows per implementation; cross-language delegation unverifiable &
Any WASM host serves the identical law; receipted chains cross language boundaries \\
Compounding interface &
Tuning effort forks across variants &
Improvement lands in the kernel/ontology and propagates by artifact refresh with provenance tripwires \\
Legacy metric at transition &
(Steam: efficiency) (Tubes: speed) &
(Here: exploration sharpness---the compiled policy is flatter, Table~\ref{tab:flat}) \\
\bottomrule
\end{tabular}
\caption{The structural mapping. Each row of the old-regime column names a
failure class the new regime removes by construction; the final row
records that the new regime was, at transition, worse on the old
regime's headline metric---as were steam and the transistor.}
\label{tab:mapping}
\end{table}

The composition argument deserves its quantitative form. Suppose each
host-level implementation of a policy has probability $p$ of divergence
per deployment-relevant change (platform arithmetic, version drift,
silent fallback paths). An audit over $H$ hosts that must agree has
integrity $(1-p)^H$, which decays exponentially: per-implementation
governance cannot scale to many hosts, exactly as tube machines could
not scale to many components. The single-kernel law replaces $(1-p)^H$
with a single binary question---do the pinned fixtures pass---whose cost
is constant in $H$. This is the availability-curve argument of the
transistor, transposed from components to hosts.

\section{Disanalogies, stated}

The analogy fails in three ways a reviewer should press on, and we
state them rather than leave them to be found.

\textbf{Scale and adoption.} Steam and the transistor were adopted by
economies over decades; this work is one laboratory ecosystem. The claim
demonstrated is \emph{mechanism}, not \emph{history}. Whether the field
reorganizes around manufactured governance is an empirical question this
thesis does not answer.

\textbf{The binding-constraint question.} We assert that auditability,
not performance, binds trustworthy autonomy. For many deployed agent
systems today this is arguably false---their operators accept
unauditable allocation in exchange for results. The thesis's claim is
that this exchange becomes unavailable as consequence scales, which is a
prediction, not a measurement.

\textbf{The kernel is small.} The certified kernel allocates over eight
candidates with a compiled lens policy. Steam powered mills; the
transistor scaled to billions. The honest statement is that the
\emph{pattern}---certified kernel, pinned cross-host identity, deletion
of local reimplementations, receipts into process law---is demonstrated
end-to-end on a size where every element is real, and the pattern's
scaling is future work with named obstacles (lens-policy expressiveness
at scale, wasm artifact governance across supply chains, the
known-mutation limitation of the counterfactual battery).

\section{Falsifiers}

The phase-change framing would be falsified, locally, by any of:
(i)~a demonstration that per-implementation governance scales across
heterogeneous hosts without a single-kernel discipline (the composition
argument would then be unnecessary);
(ii)~divergence between the certified kernel's outputs across hosts
that survives the fixture tests (the central empirical claim would be
false);
(iii)~deployment evidence that receipted, conserving, deterministic
allocation \emph{underperforms} per-implementation tuning so severely
that operators refuse the guarantee at consequence scales where
auditability matters (the trade would be priced wrong);
or (iv)~an advance that makes verification of arbitrary floating-point
policy implementations free, removing the constraint class by other
means.

\chapter{Conclusion}
\label{ch:conclusion}

This thesis set out to answer a question disguised as a complaint: after
the adversarial audit stripped away the circular benchmarks, the trivial
theorems, the notational equivalences, and finally the reference
implementation itself, what was left? The answer is an envelope and a
law. The envelope---deterministic replay, conservation by construction,
typed refusal, receipt into process data---is what consequential
allocation must satisfy before autonomy can be delegated across trust
boundaries. The law---one certified kernel, served bit-identically to
every host through WebAssembly, with the local reimplementations deleted
rather than deprecated---is what makes the envelope cheap to verify:
constant in the number of hosts, pinned by fixtures that fail by name.

The phase-change claim was argued structurally and scoped honestly: the
same three conditions that turned the steam engine and the transistor
into regime boundaries---constraint-class removal, composition unlock,
compounding interface---hold in the small for this artifact chain, with
the characteristic signature that the new regime is initially worse on
the old regime's metric. Whether that structure compounds beyond one
laboratory is the open question the artifact chain is positioned to
answer, and the falsifiers in Chapter~\ref{ch:phase} are the terms of
its answering.

The deepest contribution may be the ledger. A discipline that deletes
its own weaker claims, with receipts, is one that can be trusted to
report its stronger ones. Governance, in the end, is not a property an
artifact has. It is a property of the process that is allowed to modify
the artifact---and this thesis has tried to be an instance of its own
thesis about it.

\appendix

\chapter{Artifact and Verification Map}
\label{app:artifacts}

All quantitative claims in the text anchor here. Repositories:
\code{seanchatmangpt/autofde-lab} (branch
\code{feat/semantic-model-manufacturing}) and the sibling
\code{ash\_autofde} project (local).

\begin{table}[h]
\centering
\small
\begin{tabular}{llp{0.5\textwidth}}
\toprule
Commit & Role & Content \\
\midrule
\commit{39d244a9} & Vendoring & bcinr-cmca crate vendored; kernel becomes default engine \\
\commit{465616a3} & WASM transport & shim crate, committed artifact, wasm-first dispatch, 7 Chicago tests \\
\commit{3e2067f9} & Receipt & verification receipt with full command/exit ladder \\
\commit{c1d04ffc} & Deletion & reference-softmax engine, tau, and its tests deleted; single-kernel law \\
\commit{11e0559} & BEAM host & \code{ash\_autofde} initial commit; wasmex engine + golden fixture \\
\commit{1617520} & BEAM update & tau\_temperature attribute removed to match plan contract \\
\bottomrule
\end{tabular}
\caption{Commit anchors for the claims in this thesis.}
\end{table}

Key artifacts: vendored kernel at \code{vendor/bcinr} (pinned
\commit{2ae60cf0}; provenance in \code{VENDORING.md}); WASM shim at
\code{wasm/bcinr-cmca-wasm} with committed artifact
\code{wasm/artifacts/bcinr\_cmca\_wasm.wasm} (421{,}988 bytes; exports
\code{cmca\_alloc}, \code{cmca\_free}, \code{cmca\_rank\_json}); Python
transport \code{src/autofde\_lab/cmca/bcinr\_wasm.py}; BEAM host
\code{lib/ash\_autofde/wasm\_cmca.ex}; projection-only allocator
\code{src/autofde\_lab/cmca/cascade.py}; golden fixture test
\code{test/ash\_autofde\_wasm\_cmca\_test.exs}; bit-identity and
law-pinning tests under \code{tests/cmca/}.

Verification status at writing: 35 tests across \code{tests/cmca} and
\code{tests/beam} pass; consumer suites (OCEL counterfactual engine,
autodev loop, GymAct/OCEL ecosystem) pass; \code{ash\_autofde}
\code{mix test} passes 10 tests and 1 doctest; lint clean on all touched
files. Standing limitation: the counterfactual battery's mutators are
generated by the codebase they test (Chapter~\ref{ch:process}).

\chapter{The Deleted Claims, For the Record}
\label{app:deleted}

For transparency, the headlines that did not survive audit, in their own
terms: ``mathematically rigorous and empirically sound'' (the
theorems were one-line properties of any softmax); ``proved and verified
equivalence'' (a change of notation); ``empirical value analysis vs.\
alternatives'' (a comparator that existed only in comments); ``optimal
collapse risk, maximal option preservation'' (assertion, not evidence).
Their deletion commits are the citations. The reader is invited to
distrust, on this basis, exactly nothing less than everything else in
this manuscript---and to check.

\begin{thebibliography}{99}
\small

\bibitem{auer2002}
P.~Auer, N.~Cesa-Bianchi, and P.~Fischer,
``Finite-time Analysis of the Multiarmed Bandit Problem,''
\emph{Machine Learning}, 47(2--4):235--256, 2002.

\bibitem{bardeen1948}
J.~Bardeen and W.~H.~Brattain,
``The Transistor, A Semi-Conductor Triode,''
\emph{Physical Review}, 74(2):230--231, 1948.

\bibitem{goldberg1991}
D.~Goldberg,
``What Every Computer Scientist Should Know About Floating-Point
Arithmetic,'' \emph{ACM Computing Surveys}, 23(1):5--48, 1991.

\bibitem{landes1969}
D.~S.~Landes, \emph{The Unbound Prometheus: Technological Change and
Industrial Development in Western Europe from 1750 to the Present}.
Cambridge University Press, 1969.

\bibitem{mackenzie2001}
D.~MacKenzie, \emph{Mechanizing Proof: Computing, Risk, and Trust}.
MIT Press, 2001.

\bibitem{mokyr1990}
J.~Mokyr, \emph{The Lever of Riches: Technological Creativity and
Economic Progress}. Oxford University Press, 1990.

\bibitem{moore1965}
G.~E.~Moore, ``Cramming More Components onto Integrated Circuits,''
\emph{Electronics}, 38(8), 1965.

\bibitem{ocel2}
A.~F.~Ghahfarokhi, G.~Park, A.~Berti, and W.~M.~P.~van~der~Aalst,
``OCEL (Object-Centric Event Log) 2.0,'' 2022.
\daggerfootnote

\bibitem{ocpq}
K.~K\"usters and W.~M.~P.~van~der~Aalst,
``Object-Centric Process Querying,'' 2025. \daggerfootnote

\bibitem{pm4py}
A.~Berti, S.~J.~van~Zelst, and W.~M.~P.~van~der~Aalst,
``Process Mining for Python (PM4Py): Bridging the Gap Between Process-
and Data Science,'' 2019. \daggerfootnote

\bibitem{rosenberg1976}
N.~Rosenberg, \emph{Perspectives on Technology}. Cambridge University
Press, 1976.

\bibitem{rosin2011}
C.~Rosin, ``Multi-Armed Bandits with Episode Context,''
\emph{Annals of Mathematics and Artificial Intelligence}, 61(3):203--230,
2011. \daggerfootnote

\bibitem{sutton2018}
R.~S.~Sutton and A.~G.~Barto, \emph{Reinforcement Learning: An
Introduction}, 2nd edition. MIT Press, 2018.

\end{thebibliography}

\begin{flushleft}
\footnotesize
\emph{Bibliographic note.} Entries marked $\dagger$ are provisional to
this manuscript draft and must be verified against library records
(volumes, pages, exact author lists) before any submission or deposit.
The artifact citations (commits, files, measurements) are exact and
reproducible from the repository at the named commits.
\end{flushleft}

\end{document}
